\section{The Greek Fathers before Dionysius}\label{sec:fathers-before}

The writers who received the faith from the apostles and their disciples
speak of the angels as part of the Church's preaching. They confess that
the angels are creatures of the one God, made through his Word, beyond
number, ordered in ranks, free, and sent to minister to the salvation of
men; that they stand over nations, churches, and persons, and carry the
prayers of the faithful before God; and that the devil and his angels were
made good and fell by their own will. Irenaeus sets the angels in the rule
of faith that the Church ``has received from the apostles and their
disciples'': Christ will come to judge, and will send ``the angels who
transgressed and became apostates'' into everlasting fire (\emph{Adv.\
haer.} I.10.1). Origen, laying out what the Church has handed down before
he begins to inquire, names the same two articles: that the devil and his
angels exist, and ``that there are certain angels of God, and certain good
influences, which are His servants in accomplishing the salvation of men''
(\emph{De principiis}, preface 6, 10). Between the letter of Clement of
Rome and the dialogues of Methodius this teaching is stated, argued from
Scripture, defended against pagans and Gnostics, and pressed into
speculation. The Catechism calls the existence of the angels ``a truth of
faith'' and appeals to ``the unanimity of Tradition'' (CCC~328); these
writers are that Tradition's first voices, and every later Father builds on
what they wrote.

English quotations of the Fathers are from the \emph{Ante-Nicene Fathers}
(Edinburgh translation, American edition of 1885--1887), cited by each
witness's own divisions.

\sectionguard
\subsection{Clement of Rome, Ignatius, and the Martyrs}

\paragraph{Clement of Rome.} The letter of the Roman Church to the
Corinthians, written in Clement's name, turns to the angels when it
exhorts the Corinthians to obedience and concord. The good servant is
prompt in well-doing, and the pattern of promptness is the host of heaven:

\begin{quote}
Let us consider the whole multitude of His angels, how they stand ever
ready to minister to His will. For the Scripture saith, ``Ten thousand
times ten thousand stood around Him, and thousands of thousands
ministered unto Him, and cried, Holy, holy, holy, [is] the Lord of
Sabaoth; the whole creation is full of His glory.'' And let us therefore,
conscientiously gathering together in harmony, cry to Him earnestly, as
with one mouth, that we may be made partakers of His great and glorious
promises. (\emph{1 Clem.}\ 34)
\end{quote}

\noindent Clement joins Daniel's myriads (Dan 7:10) to Isaias's seraphim
(Isa 6:3) and draws the lesson of the liturgy from them: the angels
minister and cry the thrice-holy with one voice, and the Church, gathered
in concord, cries ``as with one mouth.'' Clement's letter thus already sets the assembly's praise beside the
angels' Sanctus. The same letter reads Moses' song in the Greek
Bible, ``He fixed the bounds of the nations according to the number of the
angels of God. His people Jacob became the portion of the Lord'' (\emph{1
Clem.}\ 29; Deut 32:8--9). The Douay, with the Vulgate and the Hebrew,
reads that God ``appointed the bounds of people according to the number of
the children of Israel'' (Deut 32:8); the Septuagint reading that Clement quotes carries the Greek teaching on
the angels of the nations, which Clement of Alexandria and Origen develop. Clement of
Rome also carries the contrast of Hebrews: Christ is ``by so much greater
than the angels, as He hath by inheritance obtained a more excellent name
than they,'' for God ``maketh His angels spirits, and His ministers a flame
of fire,'' while to the Son he says, ``Thou art my Son'' (\emph{1 Clem.}\
36; Heb 1:4--5, 7; Ps 103[104]:4).

\paragraph{Ignatius of Antioch.} On his way to martyrdom Ignatius writes to
the Trallians that he could tell them of heavenly things, but fears to
harm them as infants in Christ: ``even I, though I am bound [for Christ],
yet am not on that account able to understand heavenly things, and the
places of the angels, and their gatherings under their respective
princes, things visible and invisible'' (\emph{Trall.}\ 5). The bishop
knows that the angels have their stations, that they are gathered into
companies, and that the companies stand under princes; he treats this as
the higher teaching of the mature, and holds himself still a learner
before it. The longer recension of the letter, an interpolated expansion
of the Ignatian text, names what Ignatius left unnamed: ``the angelic
orders, and the different sorts of angels and hosts, the distinctions
between powers and dominions, and the diversities between thrones and
authorities, the mightiness of the \AE ons, and the pre-eminence of the
cherubim and seraphim'' (\emph{Trall.}\ 5, longer recension).

Ignatius sets the whole of this heavenly order under the Cross. To the
Smyrnaeans he writes: ``Let no man deceive himself. Both the things which
are in heaven, and the glorious angels, and rulers, both visible and
invisible, if they believe not in the blood of Christ, shall, in
consequence, incur condemnation'' (\emph{Smyrn.}\ 6). No rank in heaven
stands outside the economy of the blood of Christ. The prince of this
world, on the other side, is kept in ignorance of it: ``the virginity of
Mary was hidden from the prince of this world, as was also her offspring,
and the death of the Lord; three mysteries of renown, which were wrought
in silence by God'' (\emph{Eph.}\ 19). When the star shone above all the
other stars, ``every kind of magic was destroyed, and every bond of
wickedness disappeared; ignorance was removed, and the old kingdom
abolished, God Himself being manifested in human form'' (\emph{ibid.}).
The Church's worship breaks the same power: ``when ye assemble frequently
in the same place, the powers of Satan are destroyed, and the destruction
at which he aims is prevented by the unity of your faith'' (\emph{Eph.}\
13).

\paragraph{Barnabas, the Letter to Diognetus, and Polycarp.} The Epistle
of Barnabas opens its teaching on the two ways with the angels who preside
over them: ``There are two ways of doctrine and authority, the one of
light, and the other of darkness \dots\ For over one are stationed the
light-bringing angels of God, but over the other the angels of Satan''
(\emph{Barn.}\ 18). The Letter to Diognetus measures the Incarnation by
the angels: God did not send to men ``any servant, or angel, or ruler, or
any one of those who bear sway over earthly things, or one of those to
whom the government of things in the heavens has been entrusted, but the
very Creator and Fashioner of all things'' (\emph{Diogn.}\ 7). The author
takes for granted that angels bear rule over earthly and heavenly things;
the point is that the Son is above them all. At the stake Polycarp blesses
``the God of angels and powers, and of every creature'' (\emph{Mart.\
Pol.}\ 14), and the Church of Smyrna says of the martyrs who endured
before him that the Lord revealed to them the good things laid up for
those who endure, ``inasmuch as they were no longer men, but had already
become angels'' (\emph{Mart.\ Pol.}\ 2).

\sectionguard
\subsection{The Shepherd of Hermas}

No writing of the first Christian centuries is so full of angels as the
\emph{Shepherd}. Its revelations are given by the Church in the form of an
aged woman and then by an angel in the garb of a shepherd, and its
teaching on repentance is spoken throughout by angels.

\paragraph{The tower and the first-created angels.} In the third Vision the
Church shows Hermas ``a great tower, built upon the waters, of splendid
square stones,'' built by six young men while ``myriads of men were
carrying stones to it'' (\emph{Vis.}\ III.2). The tower is the Church
herself, ``founded on the word of the almighty and glorious Name'' (III.3).
Hermas asks who the builders are:

\begin{quote}
``These are the holy angels of God, who were first created, and to whom
the Lord handed over His whole creation, that they might increase and
build up and rule over the whole creation. By these will the building of
the tower be finished.'' ``But who are the other persons who are engaged
in carrying the stones?'' ``These also are holy angels of the Lord, but
the former six are more excellent than these.'' (\emph{Vis.}\ III.4)
\end{quote}

\noindent The Church is built by angels: a first rank, created first and
charged with the whole creation, and a multitude of lesser angels who
bring the stones. The stones themselves are the apostles, bishops,
teachers, and deacons, the martyrs, and the faithful; and the young in
faith who are being fitted into the building ``are admonished by the
angels to do good'' (III.5). The ninth Similitude returns to the tower and
gives it its centre. The multitude who build ``are all glorious angels, and
by them accordingly is the Lord surrounded,'' and the gate of the tower is
the Son of God: ``The glorious man is the Son of God, and those six
glorious angels are those who support Him on the right hand and on the
left. None of these glorious angels will enter in unto God apart from
Him'' (\emph{Sim.}\ IX.12). The first-created angels stand at the Son's right hand and left and sit
in the Lord's council, and they themselves come to the Father only through
the Son. The fifth Similitude has the same order: in the parable
of the vineyard ``the stakes are the holy angels of the Lord, who keep His
people together,'' and ``the friends and fellow-councillors are the holy
angels who were first created''; the Lord ``took, therefore, as
fellow-councillors His Son and the glorious angels''; God gave the people to his Son, ``and the
Son appointed His angels over them to keep them'' (\emph{Sim.}\ V.5--6).

\paragraph{Michael and the glorious angel.} In the eighth Similitude ``a
glorious angel of the Lord, who was very tall'' stands beside a great
willow that overshadows plains and mountains, and cuts from it rods for
all who are called by the name of the Lord; later he receives the rods
back, examines them, and crowns with palms those whose rods have put out shoots and fruit
(\emph{Sim.}\ VIII.1--2). The Shepherd explains:

\begin{quote}
This great tree that casts its shadow over plains, and mountains, and all
the earth, is the law of God that was given to the whole world; and this
law is the Son of God, proclaimed to the ends of the earth \dots\ And the
great and glorious angel Michael is he who has authority over this
people, and governs them; for this is he who gave them the law into the
hearts of believers: he accordingly superintends them to whom he gave it,
to see if they have kept the same. (\emph{Sim.}\ VIII.3)
\end{quote}

\noindent Daniel's Michael, the great prince who stands for the children
of his people (Dan 12:1), is here the prince of the new people of God,
examining the fidelity of each believer. In the seventh Similitude ``the
glorious angel'' governs the discipline of Hermas's household: it is he
who has commanded that Hermas be afflicted for a time, because his house
has sinned, until they repent (\emph{Sim.}\ VII).

\paragraph{The angel of repentance.} The Shepherd himself is an angel. He
comes to Hermas ``a man of glorious aspect, dressed like a shepherd,''
and says, ``I have been sent by a most venerable angel to dwell with you
the remaining days of your life''; Hermas calls him ``the shepherd, even
the angel of repentance'' (\emph{Vis.}\ V). He is Hermas's guardian: at
the end of the book the angel who first sent him says, ``I have delivered
you and your house to the Shepherd, that you may be protected by him,''
and tells Hermas that ``he is in great honour and dignity with God, and
that he is a president with great power, and mighty in his office. To him
alone throughout the whole world is the power of repentance assigned''
(\emph{Sim.}\ X.1). The Shepherd reports on his charge: ``he had a good
report of you to me, \dots\ and he may report well of these to me, and I
to the Lord'' (X.2). Guardianship runs upward in a chain from the man to
his angel, from the angel to the glorious angel, and from him to the Lord.
Hermas gives thanks that God ``had pity on all that call upon His name;
and sent the angel of repentance to us who sinned against Him, and
renewed our spirit'' (\emph{Sim.}\ IX.14).

\paragraph{The two angels with each man.} The sixth Mandate teaches the
discernment of spirits by the two angels who attend every man:

\begin{quote}
``There are two angels with a man---one of righteousness, and the other
of iniquity.'' \dots\ ``The angel of righteousness is gentle and modest,
meek and peaceful. When, therefore, he ascends into your heart,
forthwith he talks to you of righteousness, purity, chastity,
contentment, and of every righteous deed and glorious virtue.'' \dots\
``First, he [the angel of iniquity] is wrathful, and bitter, and foolish,
and his works are evil, and ruin the servants of God.'' (\emph{Mand.}\
VI.2)
\end{quote}

\noindent Each angel is known by the thoughts he raises in the heart;
Hermas is to trust the one and ``bid farewell'' to the other. Origen cites
this Mandate, with Barnabas's two ways, as the Church's witness that good
and evil thoughts are suggested by good and evil angels (\emph{De princ.}\
III.2.4). The prophet in the assembly is filled by an angel of the same
order: when the church prays, ``then the angel of the prophetic Spirit,
who is destined for him, fills the man'' (\emph{Mand.}\ XI).

\paragraph{The devil, the angels of punishment, and the angel over the
beasts.} Hermas's devil is fearsome only to the empty. ``Fear not the
devil; for, fearing the Lord, you will have dominion over the devil, for
there is no power in him'' (\emph{Mand.}\ VII). The angel of repentance
speaks as his master: ``fear not the devil; for there is no power in him
against you, for I will be with you, the angel of repentance, who am lord
over him. The devil has fear only, but his fear has no strength''
(\emph{Mand.}\ XII.4). ``The devil can wrestle against these, overthrow
them he cannot'' (XII.5). He goes about the servants of God as a man
tests his wine-jars, passing the full ones and looking for the empty, and
he enters those whom he finds empty of faith (XII.5). The sixth Similitude
shows two further angels: a merry shepherd in yellow who is ``the angel of
luxury and deceit,'' who wears out the souls of the servants of God, and a
stern shepherd who is ``the angel of punishment; and he belongs to the
just angels, and is appointed to punish,'' who afflicts the wandering with
losses, want, and sickness until they repent (\emph{Sim.}\ VI.2--3). The
chastising angel is a holy angel; his severity serves repentance. The
fourth Vision gives the angels a charge over the animal creation: when
Hermas meets a monstrous beast on the Campanian road and escapes it by
faith, he is told that ``the Lord has sent His angel, who has rule over the
beasts, and whose name is Thegri, and has shut up its mouth'' (\emph{Vis.}\
IV.2; cf.\ Dan 6:22 [6:23]).

The \emph{Shepherd} gives the Church her earliest full picture of angelic
guardianship: a personal angel sent to each penitent, a prince over the
people of God, holy angels charged with chastisement and with the beasts,
first-created angels who build the Church and never approach the Father
except through the Son.

\sectionguard
\subsection{Justin Martyr}

\paragraph{The host of good angels.} Justin writes to the emperors to answer
the charge of atheism. The Christians deny the gods of the poets, whom he
unmasks as demons, but not the true God:

\begin{quote}
And we confess that we are atheists, so far as gods of this sort are
concerned, but not with respect to the most true God, the Father of
righteousness and temperance and the other virtues, who is free from all
impurity. But both Him, and the Son (who came forth from Him and taught
us these things, and the host of the other good angels who follow and are
made like to Him), and the prophetic Spirit, we worship and adore,
knowing them in reason and truth. (\emph{1 Apol.}\ 6)
\end{quote}

\noindent Justin sets the ``host of the other good angels'' in his
confession immediately after the Son whom they follow and to whom they are
made like. The Son is their captain; the good angels form his retinue.
The same Apology states the rule of worship in the Lord's words: ``that we
ought to worship God alone, He thus persuaded us: `The greatest
commandment is, Thou shalt worship the Lord thy God, and Him only shall
thou serve'\,'' (\emph{1 Apol.}\ 16). The Christian honours the angels as
the host of the Son and adores God alone.

\paragraph{The angels' charge and their transgression.} In the Second
Apology Justin explains why the just suffer if God is their helper. God
made the world for man, subjected earthly things to him, ordered the
heavenly elements for the seasons, and

\begin{quote}
committed the care of men and of all things under heaven to angels whom
He appointed over them. But the angels transgressed this appointment,
and were captivated by love of women, and begat children who are those
that are called demons; and besides, they afterwards subdued the human
race to themselves, partly by magical writings, and partly by fears and
the punishments they occasioned, and partly by teaching them to offer
sacrifices, and incense, and libations. (\emph{2 Apol.}\ 5)
\end{quote}

\noindent Justin reads the sons of God of Genesis 6 as angels who had
received the care of men and betrayed it (Gen 6:2--4; see
\ref{sec:before-christ}); the demons are
their offspring, and the gods of the nations are these angels and demons
under the names they gave themselves. Augustine and Chrysostom read the
sons of God as the descendants of Seth, and Aquinas follows Augustine
(\emph{De civ.\ Dei} XV.23; Chrysostom, \emph{Hom.\ in Gen.}\ 22;
\ST{I q.~51 a.~3 ad~6}; see \ref{sec:fathers-latin}). The substance of
Justin's teaching stands in either reading: God gave the angels charge
over men, and some angels betrayed their charge and made themselves the
gods of the heathen.

The fall was an act of freedom. ``God in the beginning made the race of
angels and men with free-will,'' so that ``they will justly suffer in
eternal fire the punishment of whatever sins they have committed''
(\emph{2 Apol.}\ 7). The Dialogue with Trypho states it at the end, as
Justin's last word to his Jewish hearers: God, ``wishing men and angels to
follow His will, resolved to create them free to do righteousness;
possessing reason, that they may know by whom they are created, and
through whom they, not existing formerly, do now exist.'' If Scripture
foretells that some angels and men will be punished, ``it did so because it
foreknew that they would be unchangeably [wicked], but not because God had
created them so'' (\emph{Dial.}\ 141). When Trypho protests that it is
blasphemy to say that angels sinned, Justin answers from Scripture: the
Septuagint Isaias says that ``the princes in Tanis are evil angels,'' the
devil stands among the angels before the Lord in Job and at the right hand
of the priest in Zacharias, and David sang that ``the gods of the nations
are demons'' (\emph{Dial.}\ 79; Isa 30:4, LXX; Job 1:6; Zach 3:1; Ps
95[96]:5). The devil is ``an apostate from the will of God,'' overthrown by
Christ when he asked ``to be worshipped as God, contrary to the Scripture''
(\emph{Dial.}\ 125), and Christ foretold that he ``would be sent into the
fire with his host, and the men who follow him'' (\emph{1 Apol.}\ 28).

\paragraph{The Angel who is God.} Justin's longest treatment of angels
concerns an Angel who is more than an angel. Against Trypho he argues that
the One who appeared to Abraham at Mamre, to Jacob, and to Moses in the
bush is ``another God and Lord subject to the Maker of all things; who is
also called an Angel, because He announces to men whatsoever the Maker of
all things \dots\ wishes to announce to them'' (\emph{Dial.}\ 56). He
works through Genesis 18--19 with his hearers. Of the three men at Mamre,
two were angels sent to Sodom; the third, who promised Sarah a son and
returned when Isaac was born, is called God. When the two went on to
Sodom, ``He remained behind, and communed with Abraham,'' and it is he who
``rained on Sodom sulphur and fire from the Lord out of heaven'' (Gen
19:24): ``one of the three, who is both God and Lord, and ministers to Him
who is in the heavens, is Lord of the two angels'' (\emph{Dial.}\ 56).
Trypho grants that they ``were wrong in believing that the
three who were in the tent with Abraham were all angels'' (\emph{ibid.}).
At the bush the text says that ``the Angel of the Lord appeared to him in a
flame of fire'' and then that ``the Lord called to him out of the bush''
(Exod 3:2--4), and Justin concludes that ``this same One alone who is
called an Angel, and who is God, appeared to and communed with Moses''
(\emph{Dial.}\ 60). The First Apology sums it up: the Word of God ``is
called Angel and Apostle; for He declares whatever we ought to know, and
is sent forth to declare whatever is revealed,'' and he appeared of old
``sometimes in the form of fire, and sometimes in the likeness of angels''
(\emph{1 Apol.}\ 63).

Justin's reading of the theophanies distinguishes the created angels
sharply from this Angel. The Angel is ``begotten from the Father, by His
power and will, but not by abscission''; he is not a temporary radiance of
the Father like sunlight that returns to the sun, and neither are the
angels. Some held that the Father makes his power spring forth and return,
and ``in this way, they teach, He made the angels. But it is proved that
there are angels who always exist, and are never reduced to that form out
of which they sprang'' (\emph{Dial.}\ 128). The angels are abiding
creatures with their own being. At Mamre the two angels ``are nourished in
the heavens \dots\ even though they are not nourished by food similar to
that which mortals use,'' as the manna is called angels' food; and the
Scripture that says they ate at Abraham's table speaks as we say that fire
devours all things (\emph{Dial.}\ 57; Ps 77[78]:25). Augustine reads
Mamre otherwise: God appeared ``in three men, who it is not to be doubted
were angels, although some think that one of them was Christ,'' and
Abraham and Lot addressed the Lord in them because ``God was in them as He
was wont to be in the prophets'' (\emph{De civ.\ Dei} XVI.29). Aquinas
follows Augustine: Abraham offered food to the angels as to men ``in whom,
nevertheless, he worshipped God,'' and their eating was figurative
(\ST{I q.~51 a.~2 s.c.; a.~3 ad~5}; see \ref{sec:scripture}).

\sectionguard
\subsection{Athenagoras, Tatian, and Theophilus}

\paragraph{Athenagoras: the providence of the angels.} Athenagoras's Plea
for the Christians answers the charge of atheism with the Trinity and adds
the angels to the Christian account of the divine order:

\begin{quote}
Nor is our teaching in what relates to the divine nature confined to these
points; but we recognise also a multitude of angels and ministers, whom
God the Maker and Framer of the world distributed and appointed to their
several posts by His Logos, to occupy themselves about the elements, and
the heavens, and the world, and the things in it, and the goodly ordering
of them all. (\emph{Leg.}\ 10)
\end{quote}

\noindent The angels are appointed by God through the Logos and posted to
the elements and the heavens. Athenagoras defines their office in a
sentence that governs Greek angelology after him:

\begin{quote}
For this is the office of the angels,---to exercise providence for God
over the things created and ordered by Him; so that God may have the
universal and general providence of the whole, while the particular parts
are provided for by the angels appointed over them. (\emph{Leg.}\ 24)
\end{quote}

\noindent Like men, the angels are free. ``Some, free agents, you will
observe, such as they were created by God, continued in those things for
which God had made and over which He had ordained them; but some outraged
both the constitution of their nature and the government entrusted to
them.'' Among these is ``the spirit which is about matter, who was created
by God, just as the other angels were created by Him, and entrusted with
the control of matter and the forms of matter''; he ``became negligent and
wicked in the management of the things entrusted to him,'' and others of
the angels about the first firmament ``fell into impure love of virgins,''
from whom the giants were begotten (\emph{Leg.}\ 24); Athenagoras reads
Genesis 6 with Justin. The devil is not a second
principle: nothing is really opposed to God, ``for even if anything had
placed itself in opposition to God, it would have ceased to exist''; he is
opposed only to ``the good that is in God'' (\emph{ibid.}). The fallen
angels ``haunt the air and the earth, and are no longer able to rise to
heavenly things'' (\emph{Leg.}\ 25). The demons draw men to idols, ``eager
for the blood of the sacrifices,'' and when souls are misled by empty
images the demons ``cause to flow into the mind empty visions as if coming
from the idols and the statues'' and claim for themselves the glory of any
healing (\emph{Leg.}\ 26--27). Aquinas receives the teaching that the
angels are set over corporeal things as the common teaching of the holy
doctors, quoting Augustine, Damascene, and Origen (\ST{I q.~110 a.~1}).

\paragraph{Tatian: the Logos framer of the angels.} Tatian, Justin's
hearer at Rome, who after his master's death fell away to found the
sect of the Encratites, wrote in his Address to the Greeks that ``the
Logos, too, before the creation of men, was the Framer of angels,'' and
that angels and men were each ``made free to act as it pleased, not having
the nature of good, which again is with God alone'' (\emph{Or.}\ 7). The
fall began among the angels. When men attached themselves ``to one who was
more subtle than the rest, having regard to his being the first-born, and
declared him to be God,'' the Logos excluded ``the beginner of the folly
and his adherents from all fellowship with Himself''; ``that first-begotten
one through his transgression and ignorance becomes a demon; and they who
imitated him \dots\ are become a host of demons'' (\emph{ibid.}). The
translators' note gives the Greek here as \gk{angelos pr\=otogonos}, the
first-born angel. The demons ``do not indeed die easily, for they do not
partake of flesh; but while living they practice the ways of death''; their
immortality will serve their punishment (\emph{Or.}\ 14). ``None of the
demons possess flesh; their structure is spiritual, like that of fire or
air,'' and they are seen only by those whom the Spirit of God indwells;
``on this account the nature of the demons has no place for repentance''
(\emph{Or.}\ 15); ``being smitten by the word of God, they depart in
terror'' (\emph{Or.}\ 16). Tatian says that the demons, ``having received their structure from
matter and obtained the spirit which inheres in it, became intemperate and
greedy'' (\emph{Or.}\ 12); Aquinas
treats the view that created spirits have bodies naturally united to them
and rejects it (\ST{I q.~51 a.~1}).

\paragraph{Theophilus of Antioch.} Theophilus names the devil in his account
of Eden. Eve, deceived by the serpent, became the author of sin, and behind the serpent stood ``the wicked demon, who also is called
Satan, who then spoke to her through the serpent, and who works even to
this day in those men that are possessed by him''; ``he is called `demon' and `dragon,' on account
of his [\gk{apodedrakenai}] revolting from God. For at first he was an
angel'' (\emph{Ad Autol.}\ II.28). The Greek pun joins \gk{drak\=on},
dragon, to the verb of running away: the dragon is the deserter. Satan's
envy follows the first sin into the second generation: when he saw Adam
and Eve living and begetting children, ``being carried away with spite
because he had not succeeded in putting them to death,'' and saw Abel
pleasing to God, ``he wrought upon the heart of his brother called Cain,
and caused him to kill his brother Abel. And thus did death get a
beginning in this world'' (II.29). The pagan poets spoke from the same
spirits, ``and these spirits of error themselves confess that they are
demons who also formerly inspired these writers'' when the possessed are
exorcised ``in the name of the living and true God'' (II.8).

\sectionguard
\subsection{Irenaeus of Lyons}

Irenaeus writes against the Gnostic systems in which angels or a
demiurge made the world against the will, or without the knowledge, of the
highest God. His teaching on the angels is the Church's answer: one God
made all things through his Word, and the angels are among the things made.

\paragraph{Creatures, not creators.} ``Those \dots\ who say that the world
was formed by angels \dots\ contrary to the will of Him who is the
Supreme Father, err first of all in this very point.'' If the angels acted
against God's will, they would be stronger than he; if with his will, the
work is his, ``for if He is the Former of the world, He too made the
angels, or at least was the cause of their creation'' (\emph{Adv.\ haer.}\
II.2.1--3). God does not need instruments: he ``created and made all things
by His Word, while He neither required angels to assist Him in the
production of those things which are made, nor of any power greatly
inferior to Himself.'' He gave each creature its nature, and ``conferred on
spiritual things a spiritual and invisible nature, on super-celestial
things a celestial, on angels an angelical'' (II.2.4). Man too was made by
God's own hands: ``It was not angels, therefore, who made us, nor who formed
us, neither had angels power to make an image of God \dots\ For with Him
were always present the Word and Wisdom, the Son and the Spirit, by whom
and in whom, freely and spontaneously, He made all things'' (IV.20.1). The
Father, ``being in no want of angels,'' nevertheless ``has a vast and
unspeakable number of servants. For His offspring and His similitude do
minister to Him in every respect; that is, the Son and the Holy Spirit, the
Word and Wisdom; whom all the angels serve, and to whom they are subject''
(IV.7.4).

The angels are therefore works of the Creator whom the Gnostics despise.
Irenaeus turns their own boast against them: if the things above the
heavens came forth from their spiritual seed, ``then let them inform us
what is the nature of things invisible, recount the number of the Angels,
and the ranks of the Archangels, reveal the mysteries of the Thrones, and
teach us the differences between the Dominations, Principalities, Powers,
and Virtues. But they can say nothing respecting them; therefore these
beings were not made by them'' (II.30.6). What heretic can show a work to rival
``those existences which are above heaven, and which do not pass away,
such as Angels, Archangels, Thrones, Dominions, and Powers innumerable''
(II.30.3)? Irenaeus names the orders of Scripture as he finds them in
Paul, sets no fixed rank among them, and declares them all creatures. John
proves it: ``all things, whether Angels, or Archangels, or Thrones, or
Dominions, were both established and created by Him who is God over all,
through His Word,'' for ``all things were made by Him'' (III.8.3; John
1:3; Ps 32[33]:6, 9). Since they have received a beginning, the name of
God belongs to none of them, ``so that He indeed who made all things can
alone, together with His Word, properly be termed God and Lord'' (III.8.3).
The Son, co-existing with the Father from the beginning, ``always reveals
the Father to Angels, Archangels, Powers, Virtues, and all to whom He wills
that God should be revealed'' (II.30.9).

\paragraph{The Church's prayer.} Irenaeus contrasts the Church's works of
power with the magic of the heretics: ``Nor does she perform anything by
means of angelic invocations, or by incantations, or by any other wicked
curious art; but, directing her prayers to the Lord, who made all things,
in a pure, sincere, and straightforward spirit, and calling upon the name
of our Lord Jesus Christ, she has been accustomed to work miracles for the
advantage of mankind'' (II.32.5). The Church's power flows from the name
of Christ and prayer to the Creator; the names of angels used as spells
belong to the sorcerers.

\paragraph{Freedom and the apostasy.} ``In man, as well as in angels, He
has placed the power of choice (for angels are rational beings), so that
those who had yielded obedience might justly possess what is good, given
indeed by God, but preserved by themselves'' (IV.37.1). The devil is no
exception. ``We do not find that the devil created anything whatsoever,
since indeed he is himself a creature of God, like the other angels.''
God ``made neither angels nor men so by nature''; those who do not obey
him ``are sons and angels of the devil, because they do the works of the
devil'' (IV.41.1--2). The eternal fire ``was
not originally prepared for man, but for him who beguiled man, and caused
him to offend---for him, I say, who is chief of the apostasy, and for those
angels who became apostates along with him'' (III.23.3; Matt 25:41).
In the parable of the tares ``we learn that this was the apostate angel and
the enemy, because he was envious of God's workmanship'' (IV.40.3; Matt
13:25).

Irenaeus gives the devil's sin a precise root. The devil was ``one among
those angels who are placed over the spirit of the air, as the Apostle Paul
has declared in his Epistle to the Ephesians''; ``becoming envious of man,
[he] was rendered an apostate from the divine law: for envy is a thing
foreign to God'' (V.24.4; Eph 2:2). His whole power ``consists in
transgression and apostasy,'' and he can go only so far as to ``deceive and
lead astray the mind of man into disobeying the commandments of God''
(V.21.3; V.24.3); in Eden ``he disputed about God, as if God was not
there, for he was ignorant of the greatness of God,'' and was proved ``a
liar and a murderer'' (V.23.1--2; John 8:44). Satan, whose Hebrew name
``signifies an apostate,'' was overcome by the Son of man with the words
of the Law in the three temptations, so that ``as in the beginning he
enticed man to transgress his Maker's law, and thereby got him into his
power,'' he should be bound by man ``with the same chains with which he had
bound man'' (V.21.2--3). This is the recapitulation: Christ ``summed up in
Himself this enmity'' of Gen 3:15, being ``made man from a woman,'' and
trod on the serpent's head (IV.40.3). The Word ``conquering him by means of
human nature, and showing him to be an apostate, has, on the contrary, put
him under the power of man'' (V.24.4; Luke 10:19). Irenaeus also holds,
with Justin, that before the Flood ``the angels that sinned had commingled''
with the race of men (IV.36.4), and that Enoch, ``although he was a man,''
was God's legate to the angels ``because the angels when they had
transgressed fell to the earth for judgment'' (IV.16.2).

\sectionguard
\subsection{Clement of Alexandria}

Clement's \emph{Stromata} present the Christian \gk{gn\=ostikos}, the
perfect believer who knows God, as one who lives in the company of the
angels and climbs the ladder of their order toward God.

\paragraph{The ordered ascent.} The seventh book sets the angel in the scale of being: ``the best thing on earth is the most pious
man; and the best thing in heaven, the nearer in place and purer, is an
angel, the partaker of the eternal and blessed life. But the nature of the
Son, which is nearest to Him who is alone the Almighty One, is the most
perfect.'' To the Son ``is placed in subjection all the host of angels and
gods'' (\emph{Strom.}\ VII.2). All government descends from him through
ranks:

\begin{quote}
But, as I think, characteristic of the highest power is the accurate
scrutiny of all the parts, reaching even to the minutest, terminating in
the first Administrator of the universe, who by the will of the Father
directs the salvation of all; some overlooking, who are set under others,
who are set over them, till you come to the great High Priest. For on one
original first Principle, which acts according to the [Father's] will, the
first and the second and the third depend. Then at the highest extremity
of the visible world is the blessed band of angels; and down to ourselves
there are ranged, some under others, those who, from One and by One, both
are saved and save. (\emph{Strom.}\ VII.2)
\end{quote}

\noindent Clement illustrates the chain with the magnet: ``As, then, the
minutest particle of steel is moved by the spirit of the Heraclean stone,
when diffused over many steel rings; so also, attracted by the Holy
Spirit, the virtuous are added by affinity to the first abode, and the
others in succession down to the last'' (\emph{ibid.}). The same power
passes down the whole series of rings, from the Son through the angels to
the least of the saints; each ring is drawn and draws. Clement has here the
principle that Dionysius will set out as hierarchy: an order in which the
higher receive from God and hand on to the lower, so that the lower are
brought up to God through them. The angels, too, had to learn: ``we also
have already heard that angels learned the truth, and their rulers over
them; for they had a beginning'' (\emph{Strom.}\ VI.7), and their teacher
is the First-begotten, by whom all things were made; God himself is
``separated as He is from even the archangels'' (\emph{ibid.}).

At the summit Clement places the seven: ``The first-born princes of the
angels, who have the greatest power, are seven'' (\emph{Strom.}\ VI.16),
the heptad that he finds throughout creation, from the planets to the
phases of the moon (cf.\ Tob 12:15; Apoc 8:2).

\paragraph{The angels of nations and of persons.} Clement joins the Song of
Moses to the government of the world: ``by an ancient and divine order the
angels are distributed among the nations. But the glory of those who
believe is `the Lord's portion'\,'' (\emph{Strom.}\ VII.2; Deut 32:8--9,
LXX). The providence of God reaches men through ministers: ``the divine
will being conveyed to human souls, particular divine ministers
contributing to such services. For regiments of angels are distributed
over the nations and cities. And, perchance, some are assigned to
individuals'' (\emph{Strom.}\ VI.17). God's benefits reach men with ``angels at
the same time lending encouragement. For by angels, whether seen or not,
the divine power bestows good things. Such was the mode adopted in the
advent of the Lord'' (\emph{ibid.}). In the Church, the presbyters serve
for improvement and the deacons for ministry, and ``in both these
ministries the angels serve God, in the management of earthly affairs''
(\emph{Strom.}\ VII.1). Clement holds that the Son ``gave philosophy to the
Greeks by means of the inferior angels'' (\emph{Strom.}\ VII.2), so that
the partial truths of the Greek sages came to them through the angels set
over their nation.

\paragraph{The Church on earth and the angelic glory.} The grades of the
Church are modelled on heaven: ``the grades here in the Church, of
bishops, presbyters, deacons, are imitations of the angelic glory, and of
that economy which, the Scriptures say, awaits those who, following the
footsteps of the apostles, have lived in perfection of righteousness''
(\emph{Strom.}\ VI.13). The perfect advance through these grades of glory,
``for glory differs from glory,'' until they grow into ``a perfect man''
(\emph{ibid.}; 1 Cor 15:41; Eph 4:13).

\paragraph{Prayer with the angels.} The man of knowledge prays with the
angels and is kept by them. ``Whence, as is right, there being only one
good God, that some good things be given from Him alone, and that some
remain, we and the angels pray. But not similarly. For it is not the same
thing to pray that the gift remain, and to endeavour to obtain it for the
first time'' (\emph{Strom.}\ VII.7). The angels pray for perseverance in a
good already possessed; men pray to receive it. The perfect Christian
``prays in the society of angels, as being already of angelic rank, and he
is never out of their holy keeping; and though he pray alone, he has the
choir of the saints standing with him'' (VII.12). His last progress is to
receive directly what the angels mediated: the Saviour ``wishes him no
longer to stand in need of help by angels, but to receive it from Himself,
having become worthy, and to have protection from Himself by obedience''
(VII.13). The warning attends the ascent: the Christian ``will pray that he
may never fall from virtue,'' for ``he knows that some of the angels,
through carelessness, were hurled to the earth, not having yet quite
reached that state of oneness'' (VII.7).

\sectionguard
\subsection{Origen}

Origen of Alexandria is the first to treat the angels in a systematic
work, and the first to press the questions that the later schools
determine: when the angels were made, whether their ranks come from
nature or from merit, whether a fallen spirit can return, how angels and
demons act on the mind, and when a guardian is given. His \emph{De
principiis} survives entire only in the Latin of Rufinus, which the
English translation renders.

\subsubsection{The Church's preaching and the rational natures}

Origen opens by setting apart what the Church has defined from what
remains to be searched out. Among the articles of the apostolic preaching
he lists that every rational soul ``has a struggle to maintain with the
devil and his angels, and opposing influences, because they strive to
burden it with sins'' (\emph{De princ.}\ preface 5); that ``regarding the
devil and his angels, and the opposing influences, the teaching of the
Church has laid down that these beings exist indeed,'' and that it is held
by most ``that the devil was an angel, and that, having become an apostate,
he induced as many of the angels as possible to fall away with himself''
(preface 6); and that ``there are certain angels of God, and certain good
influences, which are His servants in accomplishing the salvation of men''
(preface 10). Their manner of being and the time of their creation Origen
reserves for inquiry.

The first book gathers the scriptural names of the heavenly and hostile
powers. There are holy angels, ``ministering spirits, sent forth to
minister for them who shall be heirs of salvation''; Paul names ``thrones,
and dominions, and principalities, and powers''; and the Apostle's praise
of Christ above ``every name that is named, not only in this world, but
also in that which is to come'' shows that there are further orders ``which
may not be named in this world, but will be named in the world to come''
(I.5.1; Heb 1:14; Col 1:16; Eph 1:21). Against them stand the devil and his angels and the
principalities with whom we wrestle (I.5.2; Eph 6:12). The nations themselves are ``called a part of the
angels,'' for the Most High ``fixed the boundaries of the nations according
to the number of the angels of God'' (I.5.2; Deut 32:8, LXX).

\paragraph{The fall of the anointed cherub.} Origen asks whether the holy
powers are holy by nature and the evil powers evil by nature, and answers
from the prophets that no creature is either. Ezechiel's lament over the
prince of Tyre cannot be spoken of a man: ``Thou hast been the seal of a
similitude, and a crown of comeliness among the delights of paradise \dots\
From the day when thou wert created along with the cherubim, I placed thee
in the holy mount of God. Thou wert in the midst of the fiery stones: thou
wert stainless in thy days, from the day when thou wert created, until
iniquities were found in thee'' (I.5.4; Ezek 28:12--15, as Origen reads
it). ``Or what fiery stones can he imagine in the midst of which any man
could live?'' The power addressed ``was formerly holy and happy; from which state
of happiness it fell from the time that iniquity was found in it, and was
hurled to the earth, and was not such by nature and creation'' (I.5.4).
Isaias's Lucifer teaches the same: ``if, as some think, he was a nature of
darkness, how is Lucifer said to have existed before? Or how could he arise
in the morning, who had in himself nothing of the light?'' The Lord
confirms it: ``Behold, I see Satan fallen from heaven like lightning''
(Luke 10:18). ``For at one time he was light,'' and he fell and ``had his
glory turned into dust''; he is the apostate, the ``fugitive,'' of the
Septuagint Job, ``Thou wilt take with a hook the apostate dragon''
(I.5.5; Isa 14:12--14; Job 40:20 [40:25], LXX). Origen's conclusion is what the
Fourth Lateran Council defines, that the devil and the other demons were
created good by nature and became evil of themselves (DS~800): spotless
purity exists essentially only in the Father, Son, and Holy Spirit, and ``no one is pure
either by essence or nature, and \dots\ no one was by nature polluted''
(I.5.5). Aquinas reads both oracles of the devil, Isaias's ``under the
figure of the prince of Babylon'' and Ezechiel's ``in the person of the King
of Tyre,'' and from Ezechiel's cherub concludes that ``the first angel who
sinned is called, not a Seraph, but a Cherub'' (\ST{I q.~63 a.~5; a.~7
ad~1}; see \ref{sec:q63}).

\paragraph{No nature incapable of good.} Origen presses the point to the
devil himself: ``not even the devil himself was incapable of good; but
although capable of admitting good, he did not therefore also desire it,
or make any effort after virtue. For, as we are taught by those quotations
which we adduced from the prophets, there was once a time when he was
good, when he walked in the paradise of God between the cherubim''
(I.8.3). The holiness of every creature is received: ``If there is any
other nature which is holy, it possesses this property of being made holy
by the reception or inspiration of the Holy Spirit, not having it by
nature'' (\emph{ibid.}).

\subsubsection{Offices and merits}

Origen holds that the diversity of the angelic offices answers to the
merits of the spirits who hold them. The opinion is his own and he sets it
out with care. The Gnostics said that one Creator could not justly make
some spirits rulers and others subjects without cause, and inferred many
creators. Origen answers that the diversity is just because it follows
desert:

\begin{quote}
A similar method must be followed in treating of the angels; nor are we to
suppose that it is the result of accident that a particular office is
assigned to a particular angel: as to Raphael, e.g., the work of curing
and healing; to Gabriel, the conduct of wars; to Michael, the duty of
attending to the prayers and supplications of mortals. For we are not to
imagine that they obtained these offices otherwise than by their own
merits, and by the zeal and excellent qualities which they severally
displayed before this world was formed; so that afterwards in the order of
archangels, this or that office was assigned to each one, while others
deserved to be enrolled in the order of angels, and to act under this or
that archangel, or that leader or head of an order. (\emph{De princ.}\
I.8.1)
\end{quote}

\noindent The same just judgment assigns angels to churches and to
persons: ``to one angel the Church of the Ephesians was to be entrusted; to
another, that of the Smyrn\ae ans; one angel was to be Peter's, another
Paul's; and so on through every one of the little ones that are in the
Church, for such and such angels as even daily behold the face of God must
be assigned to each one of them; and there must also be some angel that
encampeth round about them that fear God'' (I.8.1; Apoc 2:1, 8; Matt
18:10; Ps 33[34]:8). In the account of the end Origen gives the ranks in a
list: of those who remained in the beginning, some ``obtained, in the
ordering and arrangement of the world, the rank of angels; others that of
influences, others of principalities, others of powers, that they may
exercise power over those who need to have power upon their head. Others,
again, received the rank of thrones, having the office of judging or ruling
\dots\ others dominion'' (I.6.2). ``The angelic office does not exist except
as a consequence of their desert; nor do `powers' exercise power except in
virtue of their moral progress'' (I.8.4). The same rule holds among the
opposing powers, who ``have obtained these degrees in evil in proportion to
their conduct'' (\emph{ibid.}).

The merit is earned ``before this world was formed.'' In the second book
Origen states the ground: God in the beginning created ``so great a number
of rational or intellectual creatures \dots\ as He foresaw would be
sufficient,'' a definite number, since ``God has arranged all things in
number and measure'' (II.9.1; Wis 11:21). Having no cause of diversity in
himself, ``He created all whom He made equal and alike,'' and their free
will ``incited each one either to progress by imitation of God, or reduced
him to failure through negligence''; the variety of the world, from the
super-celestial beings ``clothed with heavenly and resplendent bodies'' to
the lot of men at birth, follows ``from the freedom of the individual
will'' (II.9.3, 6). Because the rational natures began to be, they are ``necessarily changeable
and mutable; since whatever power was in their substance was not in it by
nature, but was the result of the goodness of their Maker'' (II.9.2).

\paragraph{The end and the return.} Origen carries the same principle to the
end of all things, and there he offers his thought ``in the style of a
disputation rather than of strict definition'' (I.6.1). Those who fell
``have not been removed irrecoverably, but have been placed under the rule
of those holy and blessed orders,'' by whose aid they may be restored
(I.6.2). Of the devil and his angels he asks ``whether any of these orders
who act under the government of the devil \dots\ will in a future world be
converted to righteousness because of their possessing the faculty of
freedom of will, or whether persistent and inveterate wickedness may be
changed by the power of habit into nature,'' and leaves the answer to the
reader (I.6.3). A passage that Jerome translated in his letter to Avitus goes further:
``angels may become men or demons, and again from the latter they may rise
to be men or angels'' (fragment after I.8).

The Church judged these opinions. The synod held at Constantinople in 543
condemned the pre-existence of souls as minds that cooled from the love of
God, and the teaching that the punishment of demons will end in their
restoration (DS~403, 411; see \ref{sec:faith}). Aquinas answers the merit
of the first creation with the goodness of the whole: God made things
unequal for the perfection of the universe, and inequality from merit
belongs to the distribution of rewards (\ST{I q.~47 a.~2 and ad~3}); he
judges erroneous Origen's opinion that every created will can turn to good
or evil, since the demons are obstinate in evil (\ST{I q.~64 a.~2}; see
\ref{sec:fall}).

\subsubsection{The opposing powers and the struggle}

The third book treats the assault of the demons with a physician's
precision. Scripture names the enemy throughout: the serpent in Eden, inspired by
the devil; the lying spirit who offered to deceive Achab, ``from his own
(free) will and choice''; Satan who moved David to number the people; the
adversary of Job, ``defeated through the patience of Job''; the tempter of
the Saviour, who entered Judas (\emph{De princ.}\ III.2.1; 3 Kgs [1 Kgs]
22:20--22; 1 Par [1 Chr] 21:1; Eph 6:11--12).

Origen will not make the devil the cause of every sin. ``The more simple
among the believers'' think that ``if there were no devil, no single human
being would go astray,'' but hunger, thirst, and the other natural desires
do not come from him; man can exceed their measure unprompted (III.2.1--2).
The demons seize the occasion of the first excess: ``we receive certain
initial elements, and, as it were, seeds of sins, from those things which
we use agreeably to nature; but when we have indulged them beyond what is
proper, and have not resisted the first movements to intemperance, then the
hostile power, seizing the occasion of this first transgression, incites
and presses us hard in every way'' (III.2.2). Aquinas makes this passage
the ground of his answer that not every sin is due to the devil's
temptation (\ST{I q.~114 a.~3}; see \ref{sec:q114}).

The combat is just, because providence matches each man to his enemy as
the presidents of the games pair wrestlers of equal age and strength, ``so
that one individual fights against one temptation of the flesh, another
against a second \dots\ one has to resist this or that hostile power,
another has to combat two or three at the same time'' (III.2.3; 1 Cor
10:13). The mind is the field of both kinds of angels. Thoughts
``sometimes proceed from ourselves, and sometimes are originated by the
opposing powers; not seldom also are they suggested by God, or by the holy
angels,'' as Raphael accompanied Tobias and the angel spoke within
Zacharias (III.2.4; Tob 5--6; Zach 1:14). Yet what is suggested to the heart is ``a
(mental) commotion only, and an incitement instigating us either to good or
evil,'' and ``it is quite within our reach, when a malignant power has begun
to incite us to evil, to cast away from us the wicked suggestions''
(\emph{ibid.}).

\paragraph{The princes of the nations.} Among the powers are the princes of
this world whose wisdom Paul will not speak (1 Cor 2:6--8). Origen reads
them as spiritual rulers of nations: ``in the holy Scriptures we find that
there are princes over individual nations; as in Daniel we read that there
was a prince of the kingdom of Persia, and another prince of the kingdom of
Gr\ae cia, who are clearly shown, by the nature of the passage, to be not
human beings, but certain powers,'' and the prince of Tyre in Ezechiel is
``a kind of spiritual power'' (III.3.2; Dan 10:13, 20). Each of these
princes has his own wisdom, which he imparts to his nation, the occult
philosophy of the Egyptians or the astrology of the Chaldeans; and when
they saw the Saviour come to destroy false wisdom, ``not knowing what was
concealed within Him, they forthwith laid a snare for Him,'' and crucified
the Lord of glory (\emph{ibid.}; Ps 2:2; 1 Cor 2:8).

\subsubsection{Prayer, worship, and the angels of the nations}

Against the pagan philosopher Celsus Origen states how Christians honour
the angels and to whom they pray. Celsus had said that if the Christians
speak of angels descending to men, they must call them gods or demons.
Origen answers:

\begin{quote}
For we indeed acknowledge that angels are ``ministering spirits,'' and we
say that ``they are sent forth to minister for them who shall be heirs of
salvation;'' and that they ascend, bearing the supplications of men, to the
purest of the heavenly places in the universe, or even to supercelestial
regions purer still; and that they come down from these, conveying to each
one, according to his deserts, something enjoined by God to be conferred by
them upon those who are to be the recipients of His benefits. Having thus
learned to call these beings ``angels'' from their employments, we find
that because they are divine they are sometimes termed ``god'' in the
sacred Scriptures, but not so that we are commanded to honour and worship
in place of God those who minister to us, and bear to us His blessings.
For every prayer, and supplication, and intercession, and thanksgiving, is
to be sent up to the Supreme God through the High Priest, who is above all
the angels, the living Word and God. (\emph{C.\ Cels.}\ V.4)
\end{quote}

\noindent The angels are named from their office; they carry prayers up and
blessings down; the prayer itself goes to God through Christ, the High
Priest above the angels. ``It is enough to secure that the holy angels of
God be propitious to us, and that they do all things on our behalf, that
our disposition of mind towards God should imitate as far as it is within
the power of human nature the example of these holy angels, who again
follow the example of their God'' (V.5). The Christian gains the favour of
the angels by imitating them. The name ``demons,'' by contrast, ``is always
applied to those wicked powers, freed from the encumbrance of a grosser
body, who lead men astray'' (V.5).

The eighth book develops the teaching. Had Celsus meant by the servants
of God ``Gabriel and Michael, and the other angels and archangels,'' and
defined what he meant by worship, Origen would have answered him on that
question; but Celsus's servants of God are the demons, and to them no
honour is due; ``we worship with all our
power the one God, and His only Son, the Word and the Image of God, by
prayers and supplications; and we offer our petitions to the God of the
universe through His only-begotten Son'' (VIII.13). ``All angels are not
said to be angels of God, but only those that are blessed: those that have
fallen away into sin are called `angels of the devil'\,'' (VIII.25). If
Christians wish for friends beyond God and his Son,

\begin{quote}
we are taught that ``thousand thousands stood before Him, and ten
thousand times ten thousand ministered unto Him.'' And these, regarding
all as their relations and friends who imitate their piety towards God,
and in prayer call upon Him with sincerity, work along with them for their
salvation, appear unto them, deem it their office and duty to attend to
them, and as if by common agreement they visit with all manner of kindness
and deliverance those who pray to God, to whom they themselves also pray.
(\emph{C.\ Cels.}\ VIII.34)
\end{quote}

\noindent ``Let the learned Greeks say that the human soul at its birth is
placed under the charge of demons: Jesus has taught us not to despise even
the little ones in His Church, saying, `Their angels do always behold the
face of My Father which is in heaven'\,'' (VIII.34; Matt 18:10). The
guardian angel prays with his charge:

\begin{quote}
And the Christian will suffer nothing, for ``the angel of the Lord will
encamp about them that fear Him, and will deliver them,'' and his
``angel,'' who ``always beholds the face of his Father in heaven,'' offers
up his prayers through the one High Priest to the God of all, and also
joins his own prayers with those of the man who is committed to his
keeping. (\emph{C.\ Cels.}\ VIII.36)
\end{quote}

\noindent God ``sets His divine angels to watch over those who are worthy of
such guardianship, so that they can suffer nothing from demons,'' and the
believer who has accepted the guidance of Jesus, ``the `Angel of the great
counsel,'\,'' has nothing to fear from the whole host of demons (VIII.27;
Isa 9:6, LXX). Origen asks of God ``all help, and the guardianship of holy
and good angels, to defend us from the earth-spirits intent on lust, and
blood, and sacrificial odours'' (VIII.60). A little later he gives the teaching its most exultant form:

\begin{quote}
as the motion of the shadow follows that of the body which casts it, so
in like manner it follows, that when we have the favour of God, we have
also the good-will of all angels and spirits who are friends of God. For
they know who are worthy of the divine approval, and they are not only
well disposed to them, but they co-operate with them in their endeavours
to please God: they seek His favour on their behalf; with their prayers
they join their own prayers and intercessions for them. We may indeed
boldly say, that men who aspire after better things have, when they pray
to God, tens of thousands of sacred powers upon their side. (\emph{C.\
Cels.}\ VIII.64)
\end{quote}

\noindent Origen also sets angels over the fruits of the earth. Celsus
urged that the demons control the goods of life and must be propitiated;
Origen replies that ``it is not demons, but angels, who have been set over
the fruits of the earth, and over the birth of animals, it is the latter
that we praise and bless, as having been appointed by God over the things
needful for our race; yet even to them we will not give the honour which is
due to God. For this would not be pleasing to God, nor would it be any
pleasure to the angels themselves'' (VIII.57). Aquinas quotes Origen's
homily on Balaam's ass for the same teaching: ``the world has need of angels
who preside over beasts, and over the birth of animals, and trees, and
plants'' (\ST{I q.~110 a.~1}, citing \emph{Hom.\ xiv in Num.}).

\paragraph{The division of the nations.} Celsus held that the regions of
the earth were allotted to superintending spirits, so that each nation's
customs are rightly kept. Origen answers with the Song of Moses in the
Septuagint, ``He set the bounds of the people according to the number of
the angels of God; and the portion was His people Jacob'' (\emph{C.\
Cels.}\ V.29; Deut 32:8--9), read with the story of Babel (Gen 11:1--9).
He reads the history as a veiled mystery. The nations that
departed from the east and built the tower were ``handed over \dots\ to
angels of character more or less severe, and of a nature more or less
stern, until they had paid the penalty of their daring deeds; and they
were conducted by those angels, who imprinted on each his native language,
to the different parts of the earth according to their deserts'' (V.30).
Those who kept the original speech became the Lord's portion (V.31), and
their prince is ``far more powerful than the others,'' since he can take
from every nation those who turn to him and give them laws (V.32).
Christians therefore do not keep the customs set by the spirits of the
nations; ``we see that it is a
religious act to do away with the customs originally established in the
various places by means of laws of a better and more divine character,
which were enacted by Jesus'' (V.32).

\subsubsection{The angels of the little ones}

Origen's commentary on the saying ``See that ye despise not one of these
little ones'' is the Church's first extended reflection on the guardian
angel (Matt 18:10). The Douay reads: ``See that you despise not one of
these little ones: for I say to you, that their angels in heaven always
see the face of my Father who is in heaven.''

Origen distinguishes littleness and greatness of soul, which, unlike
stature of body, depend on free will and progress. The little ones are the
new-born who ``long for the reasonable milk,'' who need the nursing-fathers
and nursing-mothers of Isaias, and he asks ``whether it is their angels who
bring them in their bosom'' (\emph{Comm.\ in Matt.}\ XIII.26; 1 Pet 2:2;
Isa 49:22--23). So long as we are imperfect we need an angel, as Jacob
said, ``The angel who delivered me from all the evils''; but ``when we have
become perfected, and have passed through the stage of being subject to
nursing-fathers and nursing-mothers and guardians and stewards, we are meet
to be governed by the Lord Himself'' (XIII.26; Gen 48:16; Gal 4:1--2).
Moses, knowing himself among the great, refused the promise, ``My angel
shall go before you,'' and asked that the Lord himself go with him
(\emph{ibid.}; Exod 32:34; 33:15).

He then asks when the angels take up their charge: ``from what time `by
the laver of regeneration' \dots\ or, whether from birth they had been
appointed, according to the foreknowledge and predestination of God,''
and gathers texts for the second view (XIII.27; Ps 70[71]:6; Gal 1:15). A
further exposition would have the angel given after regeneration, so that
``during the period of unbelief they are under the angels of Satan; but,
after the regeneration, He who has redeemed us with His own blood consigns
us to a holy angel, who also, because of his purity, beholds the face of
God'' (XIII.28). To the objection from the house of Mary, where the
disciples said of Peter at the gate, ``It is his angel'' (Acts 12:15), Origen answers that ``the word of Rhoda was not necessarily a dogma,'' nor
perhaps the word of those who spoke without exact knowledge (XIII.28). Aquinas recites Origen's two opinions and follows Jerome, who
holds the guardian given at birth (\ST{I q.~113 a.~5}; see \ref{sec:q113}
and \ref{sec:disputed}).

\sectionguard
\subsection{Gregory the Wonderworker}

Gregory, Origen's pupil and afterwards bishop of Neocaesarea, delivered
at his departure from Caesarea a thanksgiving to his master. It contains
the most personal testimony of the age to a guardian angel. After giving
thanks to God through the Word, he turns to ``any of those beings who are
not seen, but yet are more godlike, and who have a special care for men'':

\begin{quote}
it shall be addressed to that being who, by some momentous decision, had
me allotted to him from my boyhood to rule, and rear, and train,---I mean
that holy angel of God who fed me from my youth, as says the saint dear to
God, meaning thereby his own peculiar one. (\emph{Pan.\ Or.}\ 4)
\end{quote}

\noindent Gregory reads Jacob's blessing as a confession of Jacob's own
angel (Gen 48:15--16), and applies it to himself. He cannot say whether
Jacob's angel was a created angel of high rank or ``perchance the Angel of
the Mighty Counsel Himself,'' but he praises ``that being, whosoever he is,
who has been the wonderful guide of our childhood, who in all other matters
has been in time past my beneficent tutor and guardian,'' the angel ``who
still at this present time sustains, and instructs, and conducts me''
(\emph{ibid.}). The office of tutor and guardian, he says, suits no human
friend, ``for we are all blind,'' but ``only him who sees beforehand all
that is for the good of our soul.'' The angel's greatest gift was to bring
him, a stranger from Pontus, to Origen. Gregory had set out for Berytus to
study law, and it was ``a certain divine companion and beneficent conductor
and guardian'' who carried him past it to Caesarea; ``that divine angel,
gave over this charge to him, and did, if I may so speak, perchance take
his rest here, not indeed under the pressure of labour or exhaustion of any
kind (for the generation of those divine ministers knows no weariness)''
(\emph{Pan.\ Or.}\ 5). Leaving his master, Gregory asks Origen to pray that
God ``may send us a good conductor, some angel to be our comrade on the
way'' (\emph{Pan.\ Or.}\ 19; cf.\ Tob 5).

\sectionguard
\subsection{Methodius of Olympus}

Methodius, bishop and martyr, wrote against Origen on the resurrection
and in praise of virginity; his teaching on the angels guards the
distinction of natures that Origen's speculation blurred.

\paragraph{The ninety-nine sheep.} In the \emph{Banquet of the Ten Virgins}
Thaleia expounds the parable of the lost sheep (Luke 15:4--7). The Word ``is
the chief Commander and Shepherd of the heavenly ones, whom all reasonable
creatures obey and attend, who tends in order and numbers the multitudes of
the blessed angels.'' Man was made to be part of this flock, created
incorrupt ``that he might honour the king and maker of all things,
responding to the shouts of the melodious angels which came from heaven.''
When he fell, the Lord ``came from heaven into (a human) life, leaving the
ranks and the armies of angels. For the mountains are to be explained by
the heavens, and the ninety and nine sheep by the principalities and powers
which the Captain and Shepherd left when He went down to seek the lost
one'' (\emph{Symp.}\ III.6). The number of the flock is not complete
without man; the angels and men together form one choir of praise under one
Shepherd. Methodius also teaches that children, even those born of
adultery, ``are committed to guardian angels,'' which would be impossible if
their existence were against the will of God (\emph{Symp.}\ II.6).

\paragraph{The distinction of the orders.} In the dialogue on the
resurrection Methodius answers the claim that the risen, being ``as the
angels,'' must shed their flesh (Matt 22:30). God, he says, ``ordained the
nature of immortal beings to be distributed not only among angels and
ministers, but also among principalities, and thrones, and powers. For the
race of angels is one, and that of principalities and powers another;
because immortal beings are not all of one order, and constitution, and
tribe, and family'':

\begin{quote}
And neither do the cherubim, departing from their own nature, assume the
form of angels; nor, again, do angels assume the form of the others. For
they cannot be anything but what they are and have been made. \dots\ For
each one among created things must remain in its own proper place, that
none may be wanting to any, but all may be full: heaven of angels, thrones
of powers, luminaries of ministers; and the more divine spots, and the
undefiled and untainted luminaries, with seraphim, who attend the Supreme
Council, and uphold the universe; and the world of men. (\emph{De res.}\
I.10)
\end{quote}

\noindent Men will be like the angels, not angels: Christ ``did not say `they
shall be angels,' but like angels, in being, for instance, crowned, as it
is written, with glory and honour; differing a little from the angels,
while near to being angels'' (I.12; Ps 8:6). Against Origen's passage of
minds from order to order, Methodius holds each order fixed in the nature
God gave it: ``if we granted that men are changed into angels, it would
follow that we say that angels also are changed into powers, and these
into one thing and the other, until our argument proceed too far for
safety'' (I.10). Aquinas holds the same distinction of nature, and
teaches that men are taken up into the angelic orders by grace, not by a
change of nature (\ST{I q.~108 a.~8}).

\paragraph{The devil over matter.} Photius's synopsis of the treatise
reports that Methodius, ``as was said also by Athenagoras,'' held that the
devil is ``a spirit, made by God, in the neighbourhood of matter, as of
course the rest of the angels are, and that he was entrusted with the
oversight of matter, and the forms of matter.'' The angels were made ``for
the care of the universe; in order that, while God exercises a perfect and
general supervision over the whole \dots\ the angels appointed for this
purpose take charge of particulars. Now the rest of them remained in the
positions for which God made and appointed them; but the devil was
insolent, and having conceived envy of us, behaved wickedly in the charge
committed to him'' (\emph{De res.}, in Photius, \emph{Bibl.}\ cod.~234, 7). Here the line
from Athenagoras to Methodius is explicit: the angels are the stewards of
particular providence, and the devil a steward who, out of envy of man,
betrayed his stewardship.

\paragraph{The dragon and the images of the angels.} Expounding the woman
clothed with the sun, Thekla reads the dragon as ``the devil, who lies in
wait to destroy the Christ-accepted mind of the baptized,'' and the stars
swept down by his tail as ``the bodies of heresies,'' called a third part
because they err concerning one of the three Persons (\emph{Symp.}\
VIII.10; Apoc 12:3--4). Aquinas takes the same verse of the angels drawn
into sin by the first angel (\ST{I q.~63 a.~8 s.c.}). A fragment of
Methodius on the resurrection, preserved by John Damascene in his defence
of the holy images, compares the honour given to the king's images with
the images of the angels: ``The images of God's angels, which are
fashioned of gold, the principalities and powers, we make to His honour
and glory'' (\emph{De res.}\ II, fragment).

\sectionguard
\subsection{The Ante-Nicene Teaching and Its Reception}

The Greek writers before Nicaea set down the teaching below at the loci
expounded above; Aquinas and the Church take it up at the places named in
the last column, and a dash stands where no single locus receives it.

\begin{fourstable}{Witness}{Locus}{Teaching}{Received at}
Clement of Rome & \emph{1 Clem.}\ 34; Dan 7:10; Isa 6:3 & The angels are innumerable, stand ready to God's will, and cry the Sanctus; the assembly cries with them & \ST{I q.~50 a.~3} (s.c.\ Dan 7:10); CCC 335\\
Clement of Rome & \emph{1 Clem.}\ 29; Deut 32:8, LXX & The nations are bounded by the number of the angels; Israel is the Lord's portion & \ST{I q.~108 a.~6}; \ST{I q.~113 a.~8}\\
Ignatius & \emph{Trall.}\ 5 & The angels have their stations and are gathered under princes & \ST{I q.~108 aa.~1--2}\\
Ignatius & \emph{Smyrn.}\ 6; \emph{Eph.}\ 19 & Heavenly powers stand under the blood of Christ; the mysteries of the Virgin and the Cross were hidden from the prince of this world & \ST{I q.~64 a.~1 ad~4}\\
Barnabas; Hermas & \emph{Barn.}\ 18; \emph{Mand.}\ VI.2 & Angels of light and angels of Satan attend the two ways; a good and an evil angel attend each man & \ST{I q.~113 a.~2}; \ST{I q.~114 a.~1}\\
Hermas & \emph{Vis.}\ III.4; \emph{Sim.}\ V.5--6; IX.12 & First-created angels build the Church; none comes to God except through the Son & CCC 331\\
Hermas & \emph{Sim.}\ VIII.3 & Michael has authority over the people of God and examines their keeping of the law & ---\\
Hermas & \emph{Vis.}\ V; \emph{Sim.}\ X.1; \emph{Mand.}\ XII.4--6 & The angel of repentance guards the penitent; the devil can wrestle but not overthrow the faithful & \ST{I q.~113 a.~1}; \ST{I q.~114 a.~1 ad~2}\\
Justin & \emph{1 Apol.}\ 6, 16 & The host of good angels follows the Son; worship is given to God alone & CCC 331\\
Justin & \emph{2 Apol.}\ 5, 7; \emph{Dial.}\ 141 & God gave angels the care of men; angels, like men, were made free, and some transgressed & \ST{I q.~59 a.~3}; \ST{I q.~113 a.~1}; Lateran IV (DS~800)\\
Justin; Athenagoras; Irenaeus & \emph{2 Apol.}\ 5; \emph{Leg.}\ 24; \emph{Adv.\ haer.}\ IV.36.4 & The sons of God of Gen 6 are angels who sinned with women & Read otherwise: \ST{I q.~51 a.~3 ad~6}, with Augustine\\
Justin & \emph{Dial.}\ 56--60, 128; \emph{1 Apol.}\ 63 & The Word appeared as the Angel at Mamre and the bush; the angels are abiding creatures, not emanations & \ST{I q.~51 a.~2 s.c.; a.~3 ad~5}\\
Athenagoras & \emph{Leg.}\ 10, 24 & The Logos posted the angels to the elements; God's general providence works through their particular care & \ST{I q.~110 a.~1}\\
Tatian & \emph{Or.}\ 7, 14--15 & The Logos framed the angels; the first-born angel fell and the demons with him; the demons have no repentance & \ST{I q.~63 a.~8}; \ST{I q.~64 a.~2}\\
Theophilus & \emph{Ad Autol.}\ II.28--29 & The devil was at first an angel; his envy moved Cain & \ST{I q.~63 a.~2}\\
Irenaeus & \emph{Adv.\ haer.}\ II.2; II.30; III.8.3; IV.20.1 & All orders are creatures of the one God through the Word; angels create nothing & \ST{I q.~45 a.~5} (s.c.); \ST{I q.~61 a.~1}; Lateran IV (DS~800)\\
Irenaeus & \emph{Adv.\ haer.}\ IV.37.1; IV.41.1; V.24.4 & Angels have free choice; the devil is God's creature and fell through envy of man & \ST{I q.~63 aa.~2, 4}; Lateran IV (DS~800)\\
Irenaeus & \emph{Adv.\ haer.}\ III.23.3; IV.40.1 & The eternal fire was prepared for the devil and his angels & Constantinople 543 (DS~411); \ST{I q.~64 a.~4}\\
Clement of Alexandria & \emph{Strom.}\ VII.2 & The angels are ranked some under others from the Son down to men, and draw men up as a magnet draws rings & \ST{I q.~106 aa.~1, 3}; \ST{I q.~108 a.~1}\\
Clement of Alexandria & \emph{Strom.}\ VI.17; VII.2 & Angels are set over nations and cities, and perhaps over individuals & \ST{I q.~113 aa.~2, 4}; \ST{I q.~108 a.~6}\\
Clement of Alexandria & \emph{Strom.}\ VI.13; VII.12 & The grades of the Church imitate the angelic glory; the perfect pray in the society of angels & \ST{I q.~108 a.~8}\\
Origen & \emph{De princ.}\ preface 6, 10 & The Church preaches that the devil and his angels exist and that good angels serve man's salvation & Lateran IV (DS~800); CCC 328\\
Origen & \emph{De princ.}\ I.5.4--5; I.8.3 & Ezek 28 and Isa 14 speak of an angel fallen by his own will; no creature is good or evil by nature & \ST{I q.~63 aa.~4--5; a.~7 ad~1}\\
Origen & \emph{De princ.}\ I.6.2; I.8.1, 4; II.9.6 & Offices and ranks are assigned by merit earned before the world & Answered: \ST{I q.~47 a.~2}; \ST{I q.~62 a.~6}; Constantinople 543 (DS~403)\\
Origen & \emph{De princ.}\ I.6.3 & Whether the demons may be converted is left to the reader & Answered: \ST{I q.~64 a.~2}; Constantinople 543 (DS~411)\\
Origen & \emph{De princ.}\ III.2.2--4 & Not every sin comes from the devil; good and evil angels suggest thoughts & \ST{I q.~111 aa.~2--3}; \ST{I q.~114 a.~3}\\
Origen & \emph{De princ.}\ III.3.2; \emph{C.\ Cels.}\ V.29--32 & Princes rule the nations; the nations were allotted to angels at Babel & \ST{I q.~108 a.~6}; \ST{I q.~113 a.~8} (the prince of Persia a good angel, with Gregory)\\
Origen & \emph{C.\ Cels.}\ V.4--5; VIII.13, 34, 36, 64 & Angels bear prayers to God and join their prayers to ours; prayer goes to God through Christ the High Priest & CCC 335--336\\
Origen & \emph{C.\ Cels.}\ VIII.57 & Angels preside over the fruits of the earth and the birth of animals & \ST{I q.~110 a.~1}\\
Origen & \emph{Comm.\ in Matt.}\ XIII.26--28 & The little ones have angels who see the Father's face; whether from birth or from baptism & \ST{I q.~113 a.~5}\\
Gregory the Wonderworker & \emph{Pan.\ Or.}\ 4--5, 19 & A guardian angel rears, teaches, and conducts a man from boyhood & \ST{I q.~113 aa.~2, 5}; CCC 336\\
Methodius & \emph{Symp.}\ III.6 & The ninety-nine sheep are the angels; man completes the flock & ---\\
Methodius & \emph{De res.}\ I.10--12 & Each order keeps its nature; the risen are like the angels, not angels & \ST{I q.~108 a.~8}\\
Methodius & \emph{De res.}, in Photius, \emph{Bibl.}\ cod.~234, 7 & The devil was set over matter and fell through envy of man & \ST{I q.~110 a.~1}; \ST{I q.~63 a.~2}\\
\end{fourstable}
